\documentclass[12pt]{article}
\usepackage[T1]{fontenc}
\usepackage[utf8]{inputenc}
\usepackage{lmodern}
\usepackage{textcomp}
\usepackage{enumitem}
\usepackage{geometry}
\geometry{margin=1in}
\begin{document}
\begin{center}
Answer Key - Psychology - Undergraduate Level Assessment 18 \
\small (Answers are ordered exactly as in the question paper.)
\end{center}

\section*{Multiple Choice}
\begin{enumerate}[label=\arabic*.]
\item \textbf{Answer:} A
\\ \textit{Options:} A) Onset phase, abuse phase, dependency phase, abstinence phase; B) Prealcoholic phase, alcoholic phase, recovery phase, relapse phase; C) Initial phase, adaptive phase, dependent phase, deteriorative phase; D) Experimental phase, social use phase, intensive use phase, compulsive use phase; E) Initial phase, middle phase, crucial phase, termination phase
\\ \textit{Note:} The correct answer reflects Jellinek's model of alcoholism, which outlines a systematic progression through distinct phases-onset, abuse, dependency, and abstinence-highlighting the typical trajectory of an alcoholic's relationship with alcohol.
\item \textbf{Answer:} E
\\ \textit{Options:} A) Arithmetic and Digit Span subtests; B) All subtests except Block Design; C) Similarities and Picture Completion subtests; D) Vocabulary subtest only; E) All Performance subtests
\\ \textit{Note:} The classic aging pattern is characterized by a decline in performance on all Performance subtests of the Wechsler Adult Intelligence Scale-Revised, reflecting the impact of aging on non-verbal cognitive abilities while maintaining relatively stable verbal abilities.
\item \textbf{Answer:} B
\\ \textit{Options:} A) Motivation can exist without any form of arousal.; B) Arousal level is a central aspect of motivation.; C) Motivation is the cause of arousal.; D) Motivation is a by-product of arousal.; E) High arousal always maximizes motivation.
\\ \textit{Note:} The correct answer is supported by the fact that the arousal theory of motivation posits that an optimal level of arousal is necessary for individuals to perform effectively and remain motivated towards their goals, indicating that arousal directly influences motivational states.
\item \textbf{Answer:} A
\\ \textit{Options:} A) decreasing both within-group and between-group variability.; B) increasing both within-group and between-group variability.; C) increasing between-group variability only.; D) increasing within-group variability.; E) decreasing between-group variability only.
\\ \textit{Note:} The denominator term in the F-ratio represents the error or variability within groups, so reducing both within-group and between-group variability leads to a more precise estimate of this error, thus increasing the F-ratio and indicating a stronger effect.
\item \textbf{Answer:} B
\\ \textit{Options:} A) "the message is within his "latitude of acceptance."""; B) he is in a neutral or slightly negative mood.; C) the message is outside his "latitude of acceptance."; D) everyone else in the group disagrees with the message.; E) everyone else in the group agrees with the message.
\\ \textit{Note:} According to the elaboration likelihood model, individuals in a neutral or slightly negative mood are more likely to engage in careful and thoughtful consideration of information (the central route) rather than relying on superficial cues, as they are motivated to resolve their uncertainty or discomfort.
\end{enumerate}

\section*{True / False}
\begin{enumerate}[label=\arabic*.]
\setcounter{enumi}{5}
\item \textbf{Answer:} False
\\ \textit{Options:} A) True; B) False
\\ \textit{Note:} Behavioral therapy primarily emphasizes modifying observable behaviors through techniques such as reinforcement and conditioning, rather than exploring dreams and unconscious thoughts, which is more characteristic of psychodynamic approaches.
\item \textbf{Answer:} False
\\ \textit{Options:} A) True; B) False
\\ \textit{Note:} The statement is false because while simultaneous presentation of a conditioned stimulus and an unconditioned stimulus can lead to learning, effective conditioning typically requires the conditioned stimulus to precede the unconditioned stimulus to create a strong association.
\item \textbf{Answer:} False
\\ \textit{Options:} A) True; B) False
\\ \textit{Note:} The statement is false because Erikson's stage of autonomy vs. shame and doubt occurs roughly between ages 1 to 3 years, while Freud's phallic stage occurs from ages 3 to 6 years.
\item \textbf{Answer:} False
\\ \textit{Options:} A) True; B) False
\\ \textit{Note:} Cognitive dissonance theory posits that individuals experience discomfort when their beliefs and behaviors are inconsistent, prompting them to seek consistency not only in their own beliefs but also in their evaluations of others' behaviors.
\item \textbf{Answer:} True
\\ \textit{Options:} A) True; B) False
\\ \textit{Note:} Multiple-choice questions typically provide specific definitions or options that guide students toward correct answers, whereas fill-in questions require deeper understanding and recall of concepts without prompts.
\end{enumerate}

\section*{Long Form Response}
\begin{enumerate}[label=\arabic*.]
\setcounter{enumi}{10}
\item \textbf{Answer:} Research findings indicate that gender differences in conversation styles significantly influence communication dynamics, particularly highlighting women's tendency to ask more questions. This inquisitive nature encourages dialogue and fosters a collaborative atmosphere, allowing for deeper engagement and understanding among participants. By asking questions, women often promote a more inclusive environment where all voices are heard, contrasting with more assertive communication styles typically associated with men. Consequently, this questioning approach can lead to enhanced interpersonal relationships and facilitate better problem-solving in group settings. Overall, women's propensity to ask questions not only shapes the flow of conversation but also underscores the importance of communication styles in social interactions.
\\ \textit{Note:} correct because it emphasizes how women's propensity to ask questions fosters collaboration and inclusivity in conversations, contrasting with more assertive styles, thereby shaping the dynamics of communication.
\item \textbf{Answer:} Henry Murray's theory of motivation posits that human behavior is driven by individual needs and motives, fundamentally shaping how individuals respond to their environments. This framework significantly influenced David McClelland's Achievement Motivation model, which emphasizes the role of specific needs-such as the need for achievement, affiliation, and power-in guiding behavior. McClelland built on Murray's premise by suggesting that individuals are motivated to pursue success and mastery in areas they deem important, producing approach behaviors toward positive outcomes. Conversely, the fear of failure or social disapproval can lead to avoidance behaviors, aligning with Murray's view that affective responses to situations influence motivation. Thus, McClelland's model reflects and expands upon Murray's understanding of the interplay between individual needs and behavioral outcomes.
\\ \textit{Note:} correct because it clearly connects Murray's theory of individual needs as the driving force behind behavior to McClelland's model, which specifically categorizes those needs into achievement, affiliation, and power, illustrating how they impact motivation and behavior.
\end{enumerate}

\vfill
\noindent\textit{End of Answer Key}
\end{document}